\section{Repaso del modelo de variaciones autoencodificadoras (VAE)}

\subsection{La idea central del VAE}

El objetivo del VAE es aprender el decodificador $p_\theta(x|z)$ para maximizar la verosimilitud marginal de los datos $p_\theta(x)$. Sin embargo, calcular directamente la verosimilitud marginal requiere integrar sobre todas las posibles $z$, lo cual es inviable desde el punto de vista computacional:

$$p_\theta(x) = \int p_\theta(x|z) p(z) \, dz$$

Aunque podemos aproximar esta integral mediante estimación de Monte Carlo, esto requeriría muestrear grandes cantidades de $z$, con un costo computacional extremadamente alto.

\begin{warningbox}{La verosimilitud marginal no puede calcularse directamente}
Mediante la fórmula de Bayes $p_\theta(x) = \frac{p_\theta(x|z) p(z)}{p_\theta(z|x)}$, calcular la verosimilitud marginal requiere conocer también la distribución posterior $p_\theta(z|x)$, pero esta última es desconocida. Esto constituye un dilema de ``el huevo o la gallina'', y es precisamente la razón fundamental por la cual el VAE introduce una aproximación variacional de la posterior, denotada como $q_\phi(z|x)$.
\end{warningbox}

\subsection{ELBO: límite inferior de la evidencia}

Dado que no se puede maximizar directamente $\log p_\theta(x)$, el VAE recurre a maximizar su límite inferior, conocido como \textbf{límite inferior de la evidencia} (Evidence Lower Bound, ELBO).

La derivación del ELBO se basa en la desigualdad de Jensen. Para cualquier distribución $q_\phi(z|x)$:

$$\log p_\theta(x) = \log \int p_\theta(x|z) p(z) \, dz$$

Introduciendo $q_\phi(z|x)$ en la integral:

$$= \log \int \frac{p_\theta(x|z) p(z)}{q_\phi(z|x)} q_\phi(z|x) \, dz = \log \mathbb{E}_{q_\phi(z|x)} \left[ \frac{p_\theta(x|z) p(z)}{q_\phi(z|x)} \right]$$

Aplicando la desigualdad de Jensen (ya que $\log$ es una función cóncava):

$$\geq \mathbb{E}_{q_\phi(z|x)} \left[ \log \frac{p_\theta(x|z) p(z)}{q_\phi(z|x)} \right] = \text{ELBO}$$

\begin{importantbox}{Descomposición del ELBO}
El ELBO puede descomponerse en dos términos:
$$\text{ELBO} = \underbrace{\mathbb{E}_{q_\phi(z|x)}[\log p_\theta(x|z)]}_{\text{Término de reconstrucción}} - \underbrace{D_{\text{KL}}(q_\phi(z|x) \| p(z))}_{\text{Término de regularización KL}}$$

\begin{itemize}
    \item \textbf{Término de reconstrucción}: fomenta que el decodificador reconstruya con precisión la entrada $x$ a partir de la variable latente $z$, siendo esencialmente equivalente a la pérdida de reconstrucción de un autoencoder.
    \item \textbf{Término de regularización KL}: alienta a que la aproximación variacional de la posterior, $q_\phi(z|x)$, se acerque lo más posible a la prior $p(z) = \mathcal{N}(0, I)$.
\end{itemize}
\end{importantbox}

\subsection{Proceso de entrenamiento del VAE}

El VAE entrena simultáneamente el codificador y el decodificador:

\begin{itemize}
    \item \textbf{Codificador} (Encoder): parametriza la aproximación variacional de la posterior $q_\phi(z|x) = \mathcal{N}(\mu_\phi(x), \sigma^2_\phi(x) I)$, mapeando los datos a la media y la varianza en el espacio latente.
    \item \textbf{Decodificador} (Decoder): parametriza la verosimilitud $p_\theta(x|z) = \mathcal{N}(\mu_\theta(z), \sigma^2 I)$, mapeando las variables latentes de vuelta al espacio de los datos.
\end{itemize}

\begin{knowledgebox}{Truco de reparametrización (Reparameterization Trick)}
La operación de muestreo $z \sim q_\phi(z|x)$ en sí misma no es diferenciable, pero mediante el truco de reparametrización se puede evitar este problema:

$$z = \mu_\phi(x) + \sigma_\phi(x) \odot \epsilon, \quad \epsilon \sim \mathcal{N}(0, I)$$

De esta manera, $z$ es diferenciable con respecto a $\mu_\phi$ y $\sigma_\phi$, permitiendo que el gradiente se propague correctamente hacia los parámetros del codificador $\phi$.
\end{knowledgebox}

En la fase de generación, solo se utiliza el decodificador: se muestrea $z$ desde $p(z) = \mathcal{N}(0, I)$ y luego se genera $x$ a través del decodificador $p_\theta(x|z)$. El codificador se emplea únicamente durante la fase de entrenamiento para aprender una distribución de verosimilitud más adecuada.

\subsection{Las dos divergencias KL gaussianas en el VAE}

Durante el entrenamiento del VAE, es necesario calcular la divergencia KL entre dos distribuciones gaussianas isotrópicas. Para $p = \mathcal{N}(\mu, \sigma^2 I)$ y $q = \mathcal{N}(0, I)$:

$$D_{\text{KL}}(p \| q) = \frac{1}{2} \left( \|\mu\|^2 + \text{tr}(\sigma^2 I) - d - \log \det(\sigma^2 I) \right)$$

$$= \frac{1}{2} \sum_{i=1}^{d} \left( \mu_i^2 + \sigma_i^2 - 1 - \log \sigma_i^2 \right)$$

Donde $d$ es la dimensión de las variables latentes.

\subsection{Compactación del ELBO y limitaciones del VAE}

\begin{importantbox}{Condición de compactación del ELBO}
El ELBO es igual a la verosimilitud logarítmica real si y solo si la aproximación variacional de la posterior coincide completamente con la verdadera posterior:
$$\text{ELBO} = \log p_\theta(x) \iff q_\phi(z|x) = p_\theta(z|x)$$

La diferencia entre ambas es exactamente: $\log p_\theta(x) - \text{ELBO} = D_{\text{KL}}(q_\phi(z|x) \| p_\theta(z|x))$, y esta brecha se denomina \textbf{brecha de holgura} (tightness gap).
\end{importantbox}

\begin{warningbox}{Limitación fundamental del VAE}
El VAE modela tanto la distribución de verosimilitud $p_\theta(x|z)$ como la aproximación variacional de la posterior $q_\phi(z|x)$ como distribuciones gaussianas. Sin embargo, para distribuciones de datos complejas (como imágenes naturales), \textbf{que ambas la verosimilitud y la posterior sean gaussianas resulta matemáticamente imposible en la mayoría de los casos}. Esta simplificación excesiva provoca que las imágenes generadas por el VAE tiendan a ser borrosas.
\end{warningbox}

Esta limitación del VAE ha motivado directamente la exploración de modelos generativos más complejos, especialmente mediante la introducción de \textbf{estructuras jerárquicas} para aumentar la capacidad expresiva del modelo.

\subsection{Resumen de la sección}

El VAE elude el problema de la incalculabilidad de la verosimilitud marginal a través del ELBO; no obstante, su suposición gaussiana es demasiado simplista y limita la calidad generativa. Para mejorar la capacidad expresiva manteniendo la viabilidad computacional, el siguiente paso consiste en introducir una estructura jerárquica en el VAE.

